\documentclass[12pt,leqno]{article}

\usepackage{amsmath,amssymb,amsthm,colortbl,xcolor,nicefrac,setspace}
\usepackage[T1]{fontenc}
\usepackage[expert,seriftt,lucidasmallscale]{lucidabr}
\usepackage{mathrsfs,hyperref,color,graphicx}
\usepackage{booktabs,tabularx,array,needspace}
\hypersetup{hidelinks,hypertexnames=false,
  pdftitle={Purity of Pure and Fundamental Propositions},
  pdfsubject={An explanatory proof of Jacob Lettie's finite-premise theorem}}

% Self-contained Forbes / Phil-211 / Lucida setup.
\setlength{\textwidth}{6.5in}
\setlength{\textheight}{9in}
\setlength{\oddsidemargin}{0mm}
\setlength{\evensidemargin}{0mm}
\setlength{\topmargin}{0mm}
\setlength{\headheight}{0mm}
\setlength{\headsep}{0mm}
\setlength{\topskip}{0mm}
\setlength{\hoffset}{0in}
\setlength{\voffset}{0in}

\newcommand{\AND}{\ensuremath{\mathbin{\&}}}
\newcommand{\OR}{\ensuremath{\vee}}
\newcommand{\IFF}{\ensuremath{\leftrightarrow}}
\newcommand{\IC}{\ensuremath{\rightarrow}}
\newcommand{\NOT}{\ensuremath{\text{$\thicksim$}}}
\newcommand{\dentails}{\ensuremath{\vdash}}
\newcommand{\Pure}{\ensuremath{\mathrm{Pure}}}
\newcommand{\Fun}{\ensuremath{\mathrm{Fun}}}
\newcommand{\Uni}{\ensuremath{\mathrm{Uni}}}
\newcommand{\PPu}{\ensuremath{\mathrm{PP}_u}}
\newcommand{\HPE}{\ensuremath{\mathrm{H}+\mathrm{PE}}}
\newcommand{\W}{\ensuremath{\mathrm{W}}}

\theoremstyle{plain}
\newtheorem{theorem}{Theorem}
\newtheorem{lemma}{Lemma}
\renewcommand{\thelemma}{\Alph{lemma}}
\newtheorem*{proposition}{Proposition}
\renewcommand{\qedsymbol}{\(\blacksquare\)}
\onehalfspacing

\begin{document}

\begin{center}
{\Large Purity of Pure and Fundamental Propositions}

\smallskip
An explanatory proof of Jacob Lettie's theorem\\
11 November 2026
\end{center}

\noindent
Can purity itself be pure if there is a fundamental proposition? Under
Bacon's full Quantified Logical Necessity, Lettie's theorem answers
Goodman's question negatively. The argument needs only a small part of
the theory, and it does not assume that fundamentality is pure.

The theorem and its proof are due to Jacob Lettie, working with the
Purity of Pure team, a group of Claude Opus 5.5 sessions.\footnote{The
source notes are \emph{If a Proposition Is Fundamental, Purity Is Not Pure}
and \emph{Weaker Theories That Still Refute Purity of Pure}, dated
9 November 2026.} This note explains the construction and proves the exact
finite-premise result formalized in the Isabelle files supplied by Lettie.
The correspondence is recorded in the appendix. Those files are unchanged.

\section{From fundamentality to a universal proposition}

Propositions have type \(t\), and \(u=t\IC t\) is the type of unary
propositional operators. Write \(\pi=\Pure_u\). Thus the instance of
Purity of Pure at issue is
\[
 \PPu:=\Pure_{u\IC t}(\pi).
\]
It says that the predicate applying to pure unary operators is itself
pure. This is different from the bottom-type instance
\(\Pure_{t\IC t}(\Pure_t)\).

Throughout, \(x,y,z,s,r\) range over propositions and \(X,Y\) over operators
of type \(u\). We use the definitions
\[
 \Box p:=(p=\top),\qquad \Diamond p:=\NOT\Box\NOT p.
\]
Following Lettie, put
\begin{align}
 R(x,y)&:=\forall Y^u\bigl(\pi(Y)\IC(Y(y)\IC\Diamond Y(x))\bigr),
       \label{eq:R}\\
 \Uni(x)&:=\forall y^t\,R(x,y). \label{eq:Uni}
\end{align}
The relation \(R(x,y)\) says that \(x\) can have every pure property that
\(y\) has. Here \emph{property} means a propositional operator of type \(u\).
We call \(x\) \emph{universal} if this holds for every proposition \(y\).
Different properties may be possible at different possibilities; universality
does not require their joint possible instantiation.

This notion extracts what the proof needs from fundamentality.
Unary Recombination says
\[
 (\Pure(X)\AND\Fun(r))
   \IC\bigl(\Box X(r)\IC\forall q\,X(q)\bigr).
\]
Logical Purity says that closed logical terms are pure. Application says
that a pure operator applied to a pure argument has a pure value.

\begin{proposition}[Fundamental propositions are universal]\label{prop:fundamental}
Given Logical Purity, Application, and unary Recombination,
\(\forall r(\Fun(r)\IC\Uni(r))\).
\end{proposition}
\begin{proof}
Let \(r\) be fundamental. Fix a proposition \(y\) and a pure operator
\(Y\) such that \(Y(y)\). The operator
\(X:=\lambda z.\NOT Y(z)\) is pure: the closed logical term
\(\lambda W^u.\lambda z.\NOT W(z)\) is pure, and Application applies
it to the pure \(Y\).

If \(\Box\NOT Y(r)\), Recombination at \(X\) gives
\(\forall q\,\NOT Y(q)\), contradicting \(Y(y)\). Hence
\(\NOT\Box\NOT Y(r)\), which is \(\Diamond Y(r)\). Since \(y\) and
\(Y\) were arbitrary, \(\Uni(r)\).
\end{proof}

It therefore suffices to show that PP, together with suitable purity
principles, rules out every universal proposition. From this point onward
the proof contains no \(\Fun\).

\section{The construction and the idea of the proof}

We work over \HPE: classical higher-order logic H with Leibniz identity,
\(\beta\eta\)-conversion, and Propositional Equivalence at type \(t\).
The latter rule takes a \emph{theorem} \(A\IFF B\) to \(A=B\).
Modalized Functionality is not assumed.

The resulting base logic proves the S4 laws
\[
 \begin{aligned}
 \mathrm{K}:&\quad \Box(A\IC B)\IC(\Box A\IC\Box B),\\
 \mathrm{T}:&\quad \Box A\IC A,
 &\mathrm{4}:&\quad \Box A\IC\Box\Box A,
 \end{aligned}
\]
Boolean identities, and
\[
 \Box(p\IFF q)\IC p=q.
\]
It is also closed under Necessitation. These rules apply to base theorems,
including theorems containing the constants \(\Pure_\sigma\).
They do \emph{not} authorize necessitating consequences of additional
purity premises. We will mark the point where those premises enter.

Define
\begin{align}
 C(z)&:=\NOT R(z,\NOT z)\AND\NOT R(\NOT z,z), \label{eq:C}\\
 F(z)&:=\Box C(z), \label{eq:F}\\
 Q(z)&:=F(z)\IC z, \label{eq:Q}\\
 P(z)&:=\NOT z\AND\NOT F(z),&
 P'(z)&:=z\AND\NOT F(z). \label{eq:tests}
\end{align}
Each expression denotes a term built from logical vocabulary and
\(\pi\). In particular, \(F,Q,P,P'\) have type \(u\).

Lettie calls \(z\) \emph{firm} when \(F(z)\). Necessarily, neither \(z\)
nor its negation can have every pure property of the other. Firmness is
thus a modal separation condition, and its definition is symmetric in
\(z\) and \(\NOT z\).

The operator \(Q\) is designed to fix firm propositions. This requires an
identity, \(Q(z)=z\), not merely a biconditional; Lemma~\ref{lem:B}
will establish it. The two tests \(P\) and \(P'\) detect, respectively,
false and true propositions that are not firm. Their role is to force
\emph{every} value of \(Q\) to be firm.

For an operator \(Y\), write \(Y\circ Q:=\lambda z.\,Y(Q(z))\).
Here is the plan. The purity premises and PP will give
\begin{equation}
 \forall x\,F(Q(x)),
 \qquad
 \forall Y\bigl(\pi(Y)\IC\pi(Y\circ Q)\bigr).
 \label{eq:goals}
\end{equation}
Suppose that \(s\) is universal, and set \(e=Q(s)\).
The first fact makes \(e\) firm. Symmetry then makes \(\NOT e\) firm,
so \(Q(\NOT e)=\NOT e\). The second fact lets us apply universality of
\(s\) to the pure operator \(Y\circ Q\). It follows that every pure
property of \(\NOT e\) is possible for \(e\), giving \(R(e,\NOT e)\).
But firmness of \(e\) rules out exactly that relation.

The remaining work is to establish the two facts in \eqref{eq:goals}
without applying Necessitation above the added premises.

\section{The nine purity premises}

We state the finite assumptions explicitly. Abstracting the sole
nonlogical constant \(\pi\) from the terms above gives closed logical
terms \(M,M',M''\) such that
\begin{equation}
 M(\pi)=P,\qquad M'(\pi)=P',\qquad
 M''(\pi)=\lambda Y^u.\,Y\circ Q.
 \label{eq:abstractors}
\end{equation}
Their types are \((u\IC t)\IC u\), \((u\IC t)\IC u\), and
\((u\IC t)\IC(u\IC u)\), respectively. More explicitly, let \(F_c,Q_c\)
be the terms obtained by replacing \(\pi\) uniformly by a variable
\(c^{u\IC t}\). Then
\[
 \begin{aligned}
 M&:=\lambda c.\lambda z.(\NOT z\AND\NOT F_c(z)),\\
 M'&:=\lambda c.\lambda z.(z\AND\NOT F_c(z)),\\
 M''&:=\lambda c.\lambda Y.\lambda z.Y(Q_c(z)).
 \end{aligned}
\]
Equation \eqref{eq:abstractors} follows by \(\beta\)-conversion.

Write
\[
 \mathrm{App}(\sigma,\tau):=
 \forall X^{\sigma\IC\tau}\forall y^\sigma
 \bigl((\Pure(X)\AND\Pure(y))\IC\Pure(X(y))\bigr).
\]
We use the following nine ordinary premises:
\begin{align*}
 \mathrm{W1}&:\Pure(M),&
 \mathrm{W2}&:\Pure(M'),&
 \mathrm{W3}&:\Pure(M'').
\end{align*}
\[
 \begin{aligned}
 \mathrm{W4}&:\mathrm{App}(u\IC t,u),&
 \mathrm{W5}&:\mathrm{App}(u\IC t,u\IC u),\\
 \mathrm{W6}&:\mathrm{App}(u,t),&
 \mathrm{W7}&:\mathrm{App}(u,u),
 \end{aligned}
\]
and
\[
 \begin{aligned}
 \mathrm{W8}&:\Pure_t(\pi(P))\IC(\pi(P)\IC\Box\pi(P)),\\
 \mathrm{W9}&:\Pure_t(\pi(P'))\IC(\pi(P')\IC\Box\pi(P')).
 \end{aligned}
\]
The first three are instances of Logical Purity, the next four are
typed Application instances, and the last two are instances of
zeroary Exhaustion,
\[
 \Pure_t(a)\IC(a\IC\Box a).
\]
None of these premises is assumed necessary.

\begin{theorem}[Lettie]\label{thm:main}
Over \HPE, the premises W1--W9 and \PPu\ imply
\(\NOT\exists x\,\Uni(x)\).

Equivalently, let \W\ be \HPE\ with W1--W9 and the tenth ordinary premise
\(\mathrm{W10}:\exists x\,\Uni(x)\). Then
\[
 \W\dentails\NOT\PPu.
\]
\end{theorem}

We prove the theorem in three stages: establish the logical properties of
\(Q\), use the purity premises to obtain \eqref{eq:goals}, and derive the
contradiction from universality.

\section{The logical part of the construction}

All results in this section are theorems of \HPE. They mention \(\pi\)
but make no assumption about what is pure.

\begin{lemma}[Firmness persists]\label{lem:A}
\(F(x)\IC\Box F(x)\).
\end{lemma}
\begin{proof}
By definition \(F(x)=\Box C(x)\), so axiom 4 gives the result.
\end{proof}

\begin{lemma}[Q fixes firm propositions]\label{lem:B}
\(F(x)\IC Q(x)=x\).
\end{lemma}
\begin{proof}
The formula \(F(x)\IC(Q(x)\IFF x)\) is a tautology.
Necessitating this base theorem and using K gives
\(\Box F(x)\IC\Box(Q(x)\IFF x)\).
Assume that \(F(x)\). Lemma~\ref{lem:A} supplies \(\Box F(x)\),
so \(\Box(Q(x)\IFF x)\). The base-logical identity principle
\(\Box(p\IFF q)\IC p=q\) now gives \(Q(x)=x\).
\end{proof}

\begin{lemma}[Symmetry]\label{lem:C}
\(F(\NOT x)=F(x)\), and \(P'(z)=P(\NOT z)\).
\end{lemma}
\begin{proof}
Expand \(C(\NOT x)\) and use the Boolean identity \(\NOT\NOT x=x\).
By Leibniz's law,
\[
 C(\NOT x)
   =\NOT R(\NOT x,x)\AND\NOT R(x,\NOT x)
   =C(x).
\]
The second identity is commutativity of conjunction, a base theorem
obtained by PE. Applying \(\lambda p.\Box p\) to the identity gives
\(F(\NOT x)=F(x)\).

For the second claim, Leibniz's law and the same double-negation identity
give
\[
 P(\NOT z)
   =\NOT\NOT z\AND\NOT F(\NOT z)
   =z\AND\NOT F(z)
   =P'(z).
\]
\end{proof}

The choice \(Q(x)=F(x)\IC x\) has a useful consequence: a false value
of \(Q\) cannot be nonfirm. This is why the first test \(P\) matters.

\begin{lemma}\label{lem:D}
\(\NOT Q(x)\IC F(Q(x))\).
\end{lemma}
\begin{proof}
Assume that \(\NOT Q(x)\), that is, \(\NOT(F(x)\IC x)\).
Then \(F(x)\), and Lemma~\ref{lem:B} gives \(Q(x)=x\).
Substitution into \(F(x)\) yields \(F(Q(x))\).
\end{proof}

\begin{lemma}[The first test necessarily fails on Q's values]\label{lem:E}
\(\Box\NOT P(Q(x))\).
\end{lemma}
\begin{proof}
Suppose that \(P(Q(x))\). Then \(\NOT Q(x)\) and \(\NOT F(Q(x))\).
Lemma~\ref{lem:D} turns the first conjunct into \(F(Q(x))\), contradicting
the second. Thus \(\NOT P(Q(x))\) is a base theorem.
Necessitation gives the displayed result.
\end{proof}

The next lemma uses both tests. Its hypotheses are necessary purity
statements; the lemma itself is still a theorem of the base logic.

\begin{lemma}[Necessary purity of the tests suffices]\label{lem:F}
\[
 \bigl(\Box\pi(P)\AND\Box\pi(P')\bigr)\IC F(Q(x)).
\]
\end{lemma}
\begin{proof}
Put \(q:=Q(x)\). We first establish the unboxed implication
\begin{equation}
 \bigl(\pi(P)\AND\pi(P')\bigr)\IC C(q).
 \label{eq:unboxed}
\end{equation}
Assume that \(\pi(P)\) and \(\pi(P')\), and suppose, for a contradiction,
that \(C(q)\) is false. Thus
\[
 R(q,\NOT q)\OR R(\NOT q,q).
\]
The proposition \(q\) is not firm: \(F(q)=\Box C(q)\) would imply \(C(q)\)
by T. Since \(\NOT F(q)\), the contrapositive of Lemma~\ref{lem:D}
gives \(q\). Consequently,
\[
 P'(q),\qquad P(\NOT q).
\]
The first follows from \(q\AND\NOT F(q)\).
For the second, use \(\NOT\NOT q=q\) and
\(F(\NOT q)=F(q)\) from Lemma~\ref{lem:C}.

If \(R(q,\NOT q)\), instantiate its defining quantifier at the pure
operator \(P\). Since \(P(\NOT q)\), it gives \(\Diamond P(q)\).
If instead \(R(\NOT q,q)\), instantiate at the pure operator \(P'\).
Since \(P'(q)\), it gives \(\Diamond P'(\NOT q)\). But
\[
 P'(\NOT q)=P(\NOT\NOT q)=P(q)
\]
by Lemma~\ref{lem:C} and double negation. Again \(\Diamond P(q)\).

In both cases this contradicts \(\Box\NOT P(q)\), supplied by
Lemma~\ref{lem:E}. Discharging the temporary assumptions proves
\eqref{eq:unboxed} in \HPE. Necessitation and K therefore give
\[
 \bigl(\Box\pi(P)\AND\Box\pi(P')\bigr)\IC\Box C(q).
\]
Since \(\Box C(q)=F(Q(x))\), this is the required implication.
\end{proof}

\section{The purity argument}

To prove Theorem~\ref{thm:main}, we now use W1--W9. The key step makes
a true assertion of purity itself pure, so that Exhaustion applies.

\begin{lemma}\label{lem:G}
The premises W1--W9 and \PPu\ give
\[
 \Box\pi(P),\qquad \Box\pi(P'),\qquad
 \forall x\,F(Q(x)),\qquad
 \forall Y\bigl(\pi(Y)\IC\pi(Y\circ Q)\bigr).
\]
\end{lemma}
\begin{proof}
First consider \(P\). The term \(M\) is pure by W1, and \(\pi\) is
pure by \PPu. Application at W4, together with \(M(\pi)=P\), gives
\[
 \pi(P).
\]
Now \(\pi\) and \(P\) are both pure. Application at W6 gives
\[
 \Pure_t(\pi(P)).
\]
Thus \(\pi(P)\) is both true and pure. Exhaustion at W8 yields
\[
 \Box\pi(P).
\]
The same argument, using W2 and W9 in place of W1 and W8, yields
\(\Box\pi(P')\). Lemma~\ref{lem:F} then gives \(F(Q(x))\) for every \(x\).

For composition, W3 makes \(M''\) pure. With \PPu\ and Application
at W5, this makes \(M''(\pi)=\lambda Y.\,Y\circ Q\) pure.
If \(Y\) is pure, Application at W7 therefore makes \(Y\circ Q\) pure.
Generalizing \(Y\) proves the last conclusion.
\end{proof}

Exhaustion has supplied necessary purity without necessitating the
premises. PP supplies the purity of \(\pi\); Application lets
that purity be used a second time, to make \(\pi(P)\) itself a pure
proposition. This second use of the purity of \(\pi\) is what allows
Exhaustion to apply.

\section{The contradiction}

\begin{proof}[Proof of Theorem~\ref{thm:main}]
Assume that \PPu\ and W1--W9 hold, and suppose that \(s\) is universal.
Put \(e:=Q(s)\). Lemma~\ref{lem:G} makes \(e\) firm.
By symmetry \(\NOT e\) is firm too, so Lemma~\ref{lem:B} gives
\begin{equation}
 Q(\NOT e)=\NOT e. \label{eq:fixed-negation}
\end{equation}

We show that \(R(e,\NOT e)\). Let \(Y\) be a pure operator such that
\(Y(\NOT e)\). By Lemma~\ref{lem:G}, \(Y\circ Q\) is pure, and
\eqref{eq:fixed-negation} gives
\[
 (Y\circ Q)(\NOT e)=Y(Q(\NOT e))=Y(\NOT e).
\]
Thus \(Y\circ Q\) holds of \(\NOT e\). Since \(s\) is universal,
\(R(s,\NOT e)\). Apply its defining quantifier to \(Y\circ Q\).
We obtain \(\Diamond(Y\circ Q)(s)\), which by \(\beta\)-conversion is
\[
 \Diamond Y(Q(s)),\quad\text{that is,}\quad \Diamond Y(e).
\]
Since \(Y\) was arbitrary, this proves \(R(e,\NOT e)\).

But \(F(e)=\Box C(e)\) gives \(C(e)\) by T, and its first conjunct is
\(\NOT R(e,\NOT e)\). This is a contradiction. Therefore no universal
proposition exists under \PPu\ and W1--W9.

Adding W10 gives a contradiction. Since the premises are sentences and
Necessitation has been applied only to base theorems, the ordinary
deduction theorem discharges \PPu. Hence \(\W\dentails\NOT\PPu\).
\end{proof}

Combining this theorem with the proposition in Section~1 gives the
negative answer to Goodman's full-QLN question. Among the QLN principles,
only unary Recombination and two instances of zeroary Exhaustion are needed.
Purity of Fun,
uniqueness of the fundamental proposition, general Persistence, Modalized
Functionality, and higher-arity QLN are not used.

If zeroary instances are not included explicitly in QLN, unary Exhaustion
and a fundamental proposition suffice here. Given a true pure proposition
\(a\), the constant operator \(\lambda q.a\) is pure by Logical Purity and
Application. It holds of every proposition. Unary Exhaustion at a
fundamental \(r\) therefore gives \(\Box a\).

Without Exhaustion, the two particular Persistence premises
\(\pi(P)\IC\Box\pi(P)\) and \(\pi(P')\IC\Box\pi(P')\) would also supply
the needed boxes. Without either source of necessary purity, the
ordinary-premise Recombination-only question remains open.

\appendix
\Needspace{8\baselineskip}
\section*{Appendix: exact statement and verification correspondence}

The main theorem is \texttt{w\_not\_pp} in Lettie's Isabelle development.
It is a closed theorem of \texttt{U\_proves}, a calculus for H with PE,
modus ponens, Generalization, and Instantiation. There is no rule for
added premises. W1--W10 occur as the ten antecedents of
\[
 \HPE\dentails
 \mathrm{W1}\IC\bigl(\mathrm{W2}\IC\cdots
   \IC(\mathrm{W10}\IC\NOT\PPu)\cdots\bigr).
\]
The displayed implications associate to the right. Consequently the
formalization cannot apply PE or Necessitation to an antecedent as if it
were an axiom. The theorem \texttt{w\_not\_pp\_C} carries this same
implication into the application's CEV calculus.

\smallskip
\begin{singlespace}
{\small
\renewcommand{\arraystretch}{1.15}
\noindent
\begin{tabularx}{\textwidth}{@{}>{\raggedright\arraybackslash}p{0.43\textwidth}>{\raggedright\arraybackslash}X@{}}
\toprule
This note & Isabelle theorem or definition\\
\midrule
Fundamentality gives universality & \texttt{u\_uni\_ctx}; existential consequence:
  \texttt{u\_fun\_ExUni}\\
\(R,\Uni,F,Q,P,P'\) & \texttt{rcR}, \texttt{rcUni}, \texttt{rcF},
  \texttt{rcQ}, \texttt{rcP}, \texttt{rcPp}\\
\(M,M',M''\) & \texttt{rc\_P}, \texttt{rc\_Pp}, \texttt{rc\_W}\\
Lemma A & \texttt{nup\_lemA}\\
Lemma B & \texttt{nup\_lemB}\\
Lemma C & \texttt{nup\_lemC}, \texttt{nup\_ident}\\
Lemma D & \texttt{nup\_lemD}\\
Lemma E & \texttt{nup\_cor}\\
Lemma F & \texttt{nup\_step2}\\
Lemma G & \texttt{nup\_step3}\\
The final contradiction & \texttt{nup\_step4}\\
No Universal Proposition & \texttt{nup\_NUP\_U}, \texttt{u\_S}\\
The ten-premise refutation & \texttt{w\_not\_pp}, \texttt{w\_not\_pp\_C}\\
The full ordinary-premise package & \texttt{cfp2}, \texttt{cfp2\_thm}\\
\bottomrule
\end{tabularx}
\par
}
\end{singlespace}

\smallskip
The definitions and the W1--W10 inventory in this note follow the
Isabelle statement, using the conjunctive
\(F(z)=\Box(\NOT R(z,\NOT z)\AND\NOT R(\NOT z,z))\).
Lettie's main note and Metamath certificate instead write
\(F(z)=\Box\NOT(R(z,\NOT z)\OR R(\NOT z,z))\).
Their values at each argument are identical by PE. This does not by itself
identify the two operators under \(\lambda\)-abstraction over \HPE.
The corresponding abstractors in W1--W3 must therefore be read with the
chosen definition; this note does not silently interchange the two finite
premise lists.

Lettie's files in \path{theories/lettie_pp/} are included unmodified and with
permission. The repository records a successful
Isabelle2026 build/export and an audit of 17 principal endpoints with no
oracles, residual hypotheses, or unresolved unification pairs.
The complete mapping is in
\path{docs/LETTIE_PP_CORRESPONDENCE.md}.

The repository also records verification of Lettie's separately supplied
Metamath certificate \texttt{w-not-pp} on 11 November 2026.
Three verifiers were used: \texttt{metamath.exe}, \texttt{checkmm}, and
\texttt{mmverify.py}.
Its dependency trace uses only typing, substitution and conversion,
the propositional and quantifier rules, and PE, not Modalized
Functionality or the file's purity or QLN axioms. That inspection covered
the statements of the rules used, not the entire axiom base.
The Metamath file is not distributed in this repository.
This explanatory revision changes no Isabelle statement or proof.

\end{document}
